% year: תשע"ח
% type: מבחן
% semester: א
% moed: ב
% number: 3א
% points: 10
% categories: taylor, derivatives
% source: exams/מבחנים/תשעח/מועד ב תשעח.pdf

\begin{question}
(10 נק') תהי $f(x)=x^2\ln(1+x)$. חשבו את $f^{(11)}(0)$.

\underline{רמז}: אפשר להשתמש בנוסחאות מקלורין.
\end{question}

\begin{hint}
המקדם של $x^{11}$ בפולינום מקלורין של $f$ שווה ל-$\frac{f^{(11)}(0)}{11!}$. איזה איבר בפיתוח של $\ln(1+x)$ תורם ל-$x^{11}$ אחרי הכפלה ב-$x^2$?
\end{hint}

\begin{solution}
פיתוח מקלורין של $\ln(1+x)$:
\[ \ln(1+x)=x-\frac{x^2}{2}+\frac{x^3}{3}-\dots+(-1)^{n-1}\frac{x^n}{n}+o(x^n). \]
נכפיל ב-$x^2$ (עם $n=9$):
\[ f(x)=x^2\ln(1+x)=x^3-\frac{x^4}{2}+\frac{x^5}{3}-\dots+(-1)^{8}\frac{x^{11}}{9}+o(x^{11}). \]
זהו פולינום ממעלה $\le 11$ ועוד $o(x^{11})$, ולכן לפי יחידות פולינום טיילור הוא פולינום מקלורין מסדר $11$ של $f$ ($f$ גזירה אינסוף פעמים בסביבת $0$). המקדם של $x^{11}$ בפולינום מקלורין הוא $\frac{f^{(11)}(0)}{11!}$, ולכן
\[ \frac{f^{(11)}(0)}{11!}=\frac19\quad\Longrightarrow\quad f^{(11)}(0)=\frac{11!}{9}=\frac{39916800}{9}=4435200. \]
\[ \boxed{f^{(11)}(0)=\frac{11!}{9}=4435200} \]
\end{solution}
